\chapter{Resumen}

\section*{Panel de estadísticas y logros para el juez \aer}
\vspace{-0.5em}

Este Trabajo de Fin de Grado presenta el estudio, análisis, diseño e implementación de una página de estadísticas y logros para el juez online \aer. El objetivo principal del proyecto es mejorar la experiencia de los usuarios mediante una visualización clara, estructurada e intuitiva de su nivel de conocimiento y evolución a lo largo del tiempo en el campo de la programación competitiva. Asimismo, se pretende aportar una fuente adicional de motivación mediante la incorporación de técnicas de gamificación que incentiven la participación activa y el aprendizaje continuo.

Para la definición del sistema, se ha llevado a cabo una encuesta dirigida a estudiantes con experiencia o interés en programación competitiva, con el fin de identificar sus necesidades, preferencias y carencias percibidas en la plataforma actual. A partir de los resultados obtenidos, se ha elaborado una especificación de requisitos funcionales y no funcionales, estableciendo métricas relevantes para evaluar el progreso del usuario, tales como número de problemas resueltos, etc. 

Además, se ha diseñado un sistema de logros estructurado en distintos niveles y categorías, orientado a reconocer hitos significativos dentro de la plataforma y reforzar conductas positivas de aprendizaje. Este sistema busca no solo premiar el rendimiento, sino también fomentar la consistencia, la mejora progresiva y la exploración.

El desarrollo de la aplicación se ha realizado siguiendo una arquitectura modular, que garantiza la escalabilidad, mantenibilidad y reutilización de los componentes. Se ha prestado especial atención a la integración con la infraestructura existente de \aer.

Como resultado, se ha obtenido una herramienta que complementa el entorno de \aer, aportando valor tanto a estudiantes como a docentes.

\section*{Palabras clave}
\vspace{-0.5em}

\noindent juez, programación, estadísticas, gamificación, visualización, logros, datos, web, react, node